% CESC 410L lab report.
%
% NOTE TO WHOEVER EDITS THIS: keep this file free of references to any
% tooling, scripts or repository paths. It is a submitted document. Build and
% packaging instructions belong in reference_docs/report_guide.md, not here.
% Comment-only lines are stripped at package time as a backstop, but do not
% rely on that.
%
% DELIBERATE EXCEPTION: a lab report keeps its OWN preamble and does not
% \input the shared coursework preamble the homework solutions use. Three
% reasons, in order of weight:
%   1. A handed-in .tex has to compile wherever the reader puts it. A shared
%      preamble is reached by a relative path that resolves only at one exact
%      depth inside this repository, so the submitted copy would not build.
%   2. The packaging check rejects a repository path on a non-comment line,
%      and an \input of the course macros file is exactly that.
%   3. A report is prose plus figures under a title block. It wants no running
%      head, no answer box and no problem header -- which is most of what the
%      shared preamble exists to provide. It also needs listings, which the
%      shared preamble does not load.
% The overlap is four package lines, a margin and a captionsetup. That is the
% whole price, and it buys a document that stands on its own anywhere.

\documentclass[11pt]{article}

\usepackage[margin=1in]{geometry}
\usepackage{amsmath, amssymb}     % equations
\usepackage{graphicx}             % \includegraphics
\usepackage{booktabs}             % clean tables
\usepackage{listings}             % code snippets
\usepackage{xcolor}
\usepackage[colorlinks=true, allcolors=blue]{hyperref}
\usepackage{caption}

% Look for figures beside this file first, then one level up. The second entry
% lets a copy of this document sitting in a subfolder still find figs/.
\graphicspath{{./}{../}}

% Code style: Courier New equivalent at 10pt, per the CEC 320 template's
% "code in Courier New 12, single spacing". Light background, no dark blocks --
% reports may be printed.
\lstset{
  basicstyle=\ttfamily\small,
  breaklines=true,
  frame=single,
  backgroundcolor=\color{gray!8},
  showstringspaces=false,
  captionpos=b,
  language=Python,
}

\title{CESC 410L --- Lab 1: \textit{Exploring Audio Signals in Python}}
\author{Gatlin Nelson \\ Section 01}
\date{Lab start: 09/08/2026 \\ Report date: 09/08/2026}

\begin{document}
\maketitle

%======================================================================
\section{Introduction}
%======================================================================

This lab moves from the sinusoids Lab 0 generated in memory to audio read from
a file. The task is to extend the shared \texttt{dsp26} package with a
\texttt{play-plot-audio} command that opens a WAV file, plays it, reports its
header properties, and presents it in both the time and frequency domains: the
full waveform, a zoomed window, and a spectrogram drawn on linear, logarithmic
and mel frequency axes. A second task extends the command so the zoom window
can be chosen from the command line rather than hard-coded.

The motivation is that a spectrogram makes visible what a waveform hides. A
waveform shows only how loud the signal is over time; the same ten seconds
viewed as a spectrogram shows which frequencies carry that energy and how they
move, which is the representation the rest of the course builds on.

The file supplied for the first two tasks turns out not to be a recording at
all. It is ten seconds of mono audio at 16\,kHz whose entire content is two
steady sinusoids added together, and that makes it a better teaching signal
than a recording would have been: the interference between the two is visible
directly in the zoomed waveform, and Section~\ref{sec:results} works out the
rate of it from the two frequencies alone.

%======================================================================
\section{Narrative}
%======================================================================

The handout code was transcribed as given and needed four changes, all forced
by running the lab from a terminal rather than clicking through it.

\textbf{The plotting helpers had to return their figures.} As written,
\texttt{plot\_audio\_sig\_in\_td} and \texttt{plot\_audio\_sig\_in\_fd} end in
\texttt{plt.show()} and return nothing. Under a desktop that is enough, because
the window is the artifact. Run without a display, \texttt{plt.show()} does
nothing at all, so the figure is created, never drawn, and discarded with no
handle to save it from. Returning the \texttt{Figure} makes the same code work
both ways.

\textbf{The audio player blocked the run indefinitely.} The handout calls
\texttt{play\_audio()} unconditionally. It opens a Tk window and waits inside
\texttt{mainloop()} until a human closes it, which is correct for the graded
demonstration and fatal for an unattended run. The first attempt at a guard
tested whether a display existed --- and that was wrong in an instructive way.
The machine \emph{did} have a display, so the guard passed, a window opened
with nobody watching, and the command produced no output and never exited. The
symptom is easy to misread as a hang in the signal processing.

The fix was to stop asking whether a window \emph{could} open and start asking
whether one \emph{should}. Batch runs already set the non-interactive plotting
backend, so that is the signal the guard now reads. Interactively nothing
changes: the player opens exactly as the handout describes.

\textbf{Figures are written to disk, in two formats.} An option was added to
save the five figures rather than only display them, mirroring Lab 0. Both PNG
and vector copies are written; the report needs the former.

\textbf{A default had to be chosen deliberately.} The programming task asks for
the new options to carry the same defaults as the plotting function, which are
``start at zero'' and ``run to the end of the file''. Applied literally with no
further thought, a run with no arguments makes the zoomed figure an exact copy
of the full-waveform figure, and the handout's own 0.4--1.4\,s view becomes
unreachable without typing it. The options are declared with the required
defaults; when \emph{neither} is supplied the handout's window is used, so the
default output still matches the handout while any explicit window is honoured.

One check is worth recording. Two saved figures with the same byte size are a
symptom, not a coincidence --- it is what over-plotting into a reused figure
looks like. All five outputs were confirmed to differ in both size and
checksum before being accepted.

%======================================================================
\section{Artifacts}
%======================================================================

The command definition, with the two options added by the programming task:

\begin{lstlisting}[caption={\texttt{app\_cli.py} --- the new start and end options}]
@app.command()
def play_plot_audio(
    path_name: str = typer.Option(
        "", "--path-name",
        help="Optional output path name for loading an audio file."),
    start_s: float = typer.Option(
        0.0, "--start-s",
        help="Start time in seconds for the zoomed waveform figure."),
    end_s: float = typer.Option(
        None, "--end-s",
        help="End time in seconds for the zoomed waveform figure."),
):
    """Play and plot audio signal."""
    play_plot_audio_(path_name=path_name, start_s=start_s, end_s=end_s)
\end{lstlisting}

\begin{lstlisting}[caption={\texttt{recipes.py} --- choosing the zoom window}]
# Honour whatever the caller asked for, and fall back to the handout's
# window only when neither bound was given.
if start_s == 0.0 and end_s is None:
    zoom_start, zoom_end = HANDOUT_ZOOM_S      # (0.4, 1.4)
else:
    zoom_start, zoom_end = start_s, end_s

fig1 = plot_audio_sig_in_td(audio, sample_rate)
fig2 = plot_audio_sig_in_td(audio, sample_rate, zoom_start, zoom_end)
\end{lstlisting}

\textbf{How the options work.} Each \texttt{typer.Option} maps a Python
parameter to a command-line flag: the first argument is the default used when
the flag is absent, the second is the flag's spelling, and the annotated type
drives both conversion and validation, so \texttt{--start-s 0.5} arrives as the
float \texttt{0.5} and a non-numeric value is rejected before any code runs.
The help strings become the text of \texttt{--help}. The command function is a
thin shell: it parses, then hands the values to
\texttt{play\_plot\_audio\_}, which passes them to
\texttt{plot\_audio\_sig\_in\_td} as its \texttt{start\_s} and \texttt{end\_s}
arguments. No plotting code changed --- the function already accepted an
arbitrary window, so the task was to expose it, not to implement it.

\begin{figure}[h]
  \centering
  \includegraphics[width=0.8\textwidth]{figs/lab1_audio_fig2_waveform_zoom.png}
  \caption{Waveform zoomed to the handout's 0.4--1.4\,s window. This is the
           sum of two steady sinusoids seen close up: a fast carrier of
           essentially constant frequency, wrapped in an envelope that swells
           to the full $\pm 1.0$ and collapses back to nothing at a regular
           rate. The window is exactly one second wide and holds 114 of these
           lobes. Section~\ref{sec:results} derives that number from the two
           tone frequencies. Nothing about the picture changes over the rest
           of the file --- the pattern shown here simply repeats for the full
           ten seconds.}
  \label{fig:zoom}
\end{figure}

\begin{figure}[h]
  \centering
  \includegraphics[width=0.8\textwidth]{figs/lab1_audio_fig4_spec_log.png}
  \caption{Spectrogram on a logarithmic frequency axis. The signal is two
           horizontal lines and nothing else: components at 440\,Hz and
           554\,Hz, level for the full ten seconds, with no harmonics stacked
           above them, no onset and no quiet column anywhere. The logarithmic
           axis is what makes the pair legible --- their 114\,Hz of separation
           is under 1.5\% of a linear 0--8\,kHz axis, on which both would sit
           in the bottom few percent of the plot as a single thick band. The
           fine vertical ribbing is not in the audio; it is treated in
           Section~\ref{sec:results}.}
  \label{fig:speclog}
\end{figure}

\begin{figure}[h]
  \centering
  \includegraphics[width=0.8\textwidth]{figs_music/lab1_audio_prog_zoom_0.5_1.5.png}
  \caption{The programming task: the same command with
           \texttt{start\_s}~$=0.5$ and \texttt{end\_s}~$=1.5$, an arbitrary
           window the original code could not select from the command line.
           Taken from the music clip the experimental task calls for, so its
           envelope is a dense mix that never repeats, rather than the strict
           periodicity of Figure~\ref{fig:zoom}.}
  \label{fig:prog}
\end{figure}

\begin{figure}[h]
  \centering
  \includegraphics[width=0.8\textwidth]{figs_music/lab1_audio_fig5_spec_mel.png}
  \caption{Experimental task: the mel spectrogram of a ten-second excerpt of
           commercially produced music, converted to mono at 16\,kHz. Where
           Figure~\ref{fig:speclog} is two bare lines in an otherwise empty
           plane, this one is occupied nearly everywhere: several instruments
           sound at once and their harmonics fill the range up to the dark
           band at the top edge, with percussion appearing as vertical
           transients cutting across every band.}
  \label{fig:music}
\end{figure}

%======================================================================
\section{Discussion and Results}
\label{sec:results}
%======================================================================

\subsection*{Handout questions}

\begin{enumerate}
  \item \textbf{Q.} \textit{What do the three frequency-axis formats show that
        the others do not?} \\
        \textbf{A.} All three are the same short-time Fourier transform and
        differ only in how the frequency axis is mapped. Writing the transform
        explicitly,
        \begin{equation}
          X(m, k) = \sum_{n=0}^{N-1} x[n + mH]\, w[n]\,
                    e^{-j 2 \pi k n / N},
          \label{eq:stft}
        \end{equation}
        with window length $N$, hop $H$ and window $w$. A linear axis spaces
        the bins $k$ evenly, which spreads the top octave across half the plot
        and crushes everything below a kilohertz into a sliver at the bottom.
        A logarithmic axis gives equal distance to equal frequency
        \emph{ratios}: octaves become evenly spaced and the low end is
        expanded, which is exactly why the two components of
        Figure~\ref{fig:speclog} are separable there and would not be on a
        linear axis. The same warping works against a harmonic series, whose
        members are evenly spaced on a linear axis and bunch progressively
        tighter as they climb a logarithmic one. A mel axis warps the scale to
        approximate human pitch perception, roughly linear below 1\,kHz and
        logarithmic above.

  \item \textbf{Q.} \textit{Why load the file at its own sample rate?} \\
        \textbf{A.} The loader resamples to a fixed default unless told
        otherwise. Left alone, the rate printed from the file header and the
        rate labelling the spectrogram axes would describe different signals,
        and every frequency read off the figure would be wrong by the ratio
        between them.
\end{enumerate}

\subsection*{Results}

With $N = 400$ samples and $H = 160$ at a 16\,kHz sampling rate, each frame of
Equation~\ref{eq:stft} spans
\[
  \frac{N}{f_s} = \frac{400}{16000} = 25\,\text{ms},
  \qquad
  \frac{H}{f_s} = \frac{160}{16000} = 10\,\text{ms}
\]
between successive frames. This is the standard framing for speech, and the
choice is a direct trade: the 25\,ms window sets the frequency resolution at
$f_s/N = 40$\,Hz, while the 10\,ms hop sets how sharply an onset can be located
in time. Widening the window sharpens the frequency detail and smears the
onsets; narrowing it does the reverse.

Forty hertz is what makes the supplied file legible at all, since its two
components are 114\,Hz apart and so survive as two lines rather than one. The
margin is smaller than it looks. A Hann window of $N = 400$ samples has a main
lobe $4 f_s / N = 160$\,Hz wide from null to null --- wider than the separation
being resolved --- so the two peaks only just stay distinct, and the trough
between them falls a mere 9\,dB below their tops. Halving the window would
close that trough entirely and the two tones would merge into one.

The transform is taken with a 512-point FFT over a 400-sample window, so the
window is zero-padded to the next power of two. This interpolates the spectrum
onto a finer grid without adding information: it makes the display smoother,
not the resolution better, which remains $f_s/N$.

\subsection*{Beats: the supplied file is two tones, and the zoom shows it}

The default file is not a recording. Two sinusoids fitted to it account for
essentially all of its energy, at
\[
  f_1 = 440.0\ \text{Hz}, \qquad f_2 = 554.4\ \text{Hz},
\]
equal in amplitude and constant for the entire ten seconds. Their ratio is
$f_2/f_1 = 1.2599$, which is $2^{4/12}$ to five figures --- four semitones, a
major third, an A and the C\# above it. That the amplitude never varies is why
the whole-file waveform is a featureless block: correct output, not a plotting
failure.

What makes the pair worth plotting is what their sum does. Applying the
sum-to-product identity to two unit cosines,
\begin{equation}
  \cos(2\pi f_1 t) + \cos(2\pi f_2 t)
    = 2 \cos\!\left(2\pi \tfrac{f_1 + f_2}{2} t\right)
        \cos\!\left(2\pi \tfrac{f_1 - f_2}{2} t\right),
  \label{eq:beat}
\end{equation}
which recasts a sum of two tones as a single tone multiplied by a slow one.
The fast factor sits at the mean frequency,
\[
  \frac{f_1 + f_2}{2} = \frac{440.0 + 554.4}{2} = 497.2\ \text{Hz},
\]
and that is the pitch actually heard --- a note that is present in neither
component. The slow factor, at the half-difference
$(f_2 - f_1)/2 = 57.2$\,Hz, is too low to hear as a pitch and instead acts as
the envelope of the fast one.

The rate of the visible lobes is the point where care is needed. The envelope
is the \emph{magnitude} of that slow cosine, and taking the magnitude folds
its negative half up, halving its period: the envelope reaches zero twice per
cycle of the $(f_1 - f_2)/2$ cosine. The lobes therefore arrive at the full
difference frequency,
\[
  |f_2 - f_1| = 554.4 - 440.0 = 114.4\ \text{Hz},
\]
not at the 57.2\,Hz of the factor in Equation~\ref{eq:beat}. Checked against
Figure~\ref{fig:zoom}, whose window is exactly one second wide: the envelope
passes through 114 nulls across it. This is the Lab 0 material made audible
and visible at once --- two sinusoids that were added, never multiplied,
producing an amplitude modulation that no single one of them contains.

\subsection*{The spectrogram's own sampling rate, aliasing}

Figure~\ref{fig:speclog} carries a fine vertical ribbing across its lower
bands, and it is the sampling theorem appearing in a place one does not look
for it. The spectrogram advances one frame every $H/f_s = 10$\,ms, so as a
picture of how the signal changes over time it is itself a sampled signal, at
100 frames per second. The envelope it is trying to depict repeats at
114.4\,Hz, which is above that frame rate's own Nyquist limit of 50\,Hz. It
therefore cannot be represented and is aliased down to
\[
  |114.4 - 100| = 14.4\ \text{Hz},
\]
which is the rate of the ribbing. Measured frame by frame on the transform
itself, the 500\,Hz bin swings 14\,dB at 14.3\,Hz, matching the prediction.
The ribbing is not in the audio and would disappear under a shorter hop; it is
the display's sampling rate beating against the signal, in a figure whose
purpose is to teach sampling.

\subsection*{A sampling artefact visible in the music spectrogram}

Figure~\ref{fig:music} shows a dark horizontal band along its top edge that the
figures of the supplied file do not have --- they have no energy up there to
band-limit in the first place --- and it is a direct illustration of the
sampling theorem rather than a defect. The source recording is a 44.1\,kHz
stereo file
carrying content up to about 16\,kHz. Converting it to the 16\,kHz mono the
task specifies moves the Nyquist frequency to
\[
  f_N = \frac{f_s}{2} = \frac{16000}{2} = 8000\ \text{Hz},
\]
so every component above 8\,kHz must be removed before resampling --- if it
were not, those components would alias down and appear as false low
frequencies. Measured on the converted clip, energy falls away above
approximately 7.8\,kHz, and the dark band spanning roughly 7.8--8.0\,kHz is
the transition region of the anti-aliasing filter: the price paid, in bandwidth,
for a filter that cannot be infinitely sharp.

The practical result is that the same five figures can now be produced either
interactively or unattended, and the zoom window is a parameter rather than a
constant --- so examining a different second of audio is a command-line
argument instead of an edit and a re-run.

\end{document}
